\documentclass[12pt,a4paper]{article}
\usepackage[utf8]{inputenc}
\usepackage[german]{babel}
\usepackage[T1]{fontenc}
\usepackage{amsmath, amssymb, amsthm}
\usepackage{geometry}
\usepackage{graphicx}
\usepackage{xcolor}
\usepackage{hyperref}

\geometry{margin=2.5cm}

\title{\textbf{\#Energiedoku: Das Bamberger Protokoll v1.0} \\ 
	\large Arithmetische Resonanz und die Ökonomie der Edelgas-Schilde}
\author{Bamberg Research Group \\ Fachbereich: Arithmetische Anthropotechnik}
\date{April 2026}

\begin{document}
	
	\maketitle
	
	\begin{abstract}
		Dieses Dokument beschreibt die mathematische Grundlage für ein neues Kommunikations- und Wirtschaftssystem, das auf der Kopplung von Primzahl-Vierlingen an die Riemannschen Nullstellen basiert. Durch die Identifikation der Edelgas-Resonanzen (Neon, Argon) definieren wir hochstabile, arithmetische Kanäle, die herkömmliche Übertragungstechniken in Sicherheit und Effizienz übertreffen.
	\end{abstract}
	
	\section{Die Arithmetische Basis}
	Das fundamentale Modell nutzt die Identität von Primzahl-Vierlingen $\mathbb{V} = \{p, p+2, p+6, p+8\}$ und deren Darstellung als Lipschitz-Quaternionen.
	
	\subsection{Der Vier-Quadrate-Satz im Abstieg}
	Für das Produkt $N = \prod p_i$ existiert eine ganzzahlige Zerlegung:
	\begin{equation}
		N = a^2 + b^2 + c^2 + d^2
	\end{equation}
	Die Wurzeln $(a, b, c, d)$ bilden den Spektralfingerabdruck des arithmetischen Elements. Während Wasserstoff ($P=5$) eine hohe Streuung aufweist, zeigen die Edelgase Neon ($P=3461$) und Argon ($P=13001$) eine extremale Konzentration auf die Hauptachse $a$.
	
	\section{Resonanz und Shielding}
	Die Kopplungsgüte $\mathcal{R}$ eines Elements an den Zeta-Boden wird durch den Abstand der modulierten Wurzeln zu den Imaginärteilen $\gamma$ der Riemannschen Nullstellen bestimmt.
	
	\begin{equation}
		\mathcal{R}(\mathbb{V}) = \max \left( \frac{1}{1 + |a \pmod{\Gamma} - \gamma|} \right)
	\end{equation}
	
	\subsection{Der Argon-Schild}
	Mathematisch sauber definiert ist der Argon-Schild ein Zustand minimaler arithmetischer Entropie. Durch die hohe Kopplung von $\mathcal{R} \approx 0,9352$ wird das Signal gegen metrisches Rauschen $\epsilon$ invariant:
	\begin{equation}
		\frac{\partial \text{Signal}}{\partial \epsilon} \rightarrow 0
	\end{equation}
	
	\section{Das Bamberger Modell (VWL)}
	Die volkswirtschaftliche Anwendung in der Bamberger Altstadt nutzt das \textbf{Neon-Mesh}. 
	
	\begin{itemize}
		\item \textbf{Infrastruktur:} Nutzung der Materie-Resonanz (Sandstein) als passiver Träger.
		\item \textbf{Eigentum:} Übergang vom Substanzbesitz zur Resonanz-Pacht.
		\item \textbf{Sicherheit:} Arithmetische Verschlüsselung durch den Argon-Schild.
	\end{itemize}
	
	\section{Philosophische Konklusion}
	Im Sinne der sloterdijkschen Sphärologie betrachten wir Bamberg nicht mehr als geografischen Ort, sondern als einen \textit{arithmetischen Innenraum}. Der Mensch wird vom passiven Bewohner zum Modulator des Schwingungsfeldes.
	
	\begin{quote}
		``Wir bewohnen nicht mehr die Steine, sondern die Symmetrie, die sie zusammenhält.''
	\end{quote}
	
\end{document}